\section{Comparison --- which algorithm, which denoiser, where}
\label{sec:comparison}

\paragraph{Which solver wins.}
For every (structure, noise) the solver with the best NMSE and its curated metrics:
\IfFileExists{tables/cross_comparison_matirf.tex}{\begin{center}\begin{minipage}{\linewidth}\centering
\captionof{table}{Best algorithm by NMSE on MA-TIRF --- winner per structure and noise level.}
\begin{adjustbox}{max width=\linewidth}
\begin{tabular}{lllllllll}
\toprule
dataset & noise & best (by NMSE) & params & NMSE & PSNR & Depth Error (nm) & Stack Recovery & chi2\_ratio \\
\midrule
cell & 0 & \textbf{ADMM} & kappa=0.1, mu=0.5 & 0.0318 & 28.8 & 4 & 0.636 & 5.2e-08 \\
cell & 0.02 & \textbf{ADMM-PnP} & denoiser=Gaussian, sigma=50, rho=0.03 & 0.177 & 21.3 & 9 & 0 & 0.96 \\
cell & 0.05 & \textbf{ADMM-PnP} & denoiser=Gaussian, sigma=50, rho=0.03 & 0.307 & 18.9 & 21 & 0 & 0.903 \\
fibres & 0 & \textbf{ADMM} & kappa=0.01, mu=0.5 & 0.0406 & 38.0 & 10 & 0.966 & 4.86e-06 \\
fibres & 0.02 & \textbf{PNP} & denoiser=Gaussian, sigma=25, lambda\_kz=12, iter=40 & 0.385 & 28.2 & 59 & 0.0172 & 0.527 \\
fibres & 0.05 & \textbf{PNP} & denoiser=Gaussian, sigma=25, lambda\_kz=12, iter=40 & 0.653 & 25.9 & 85 & 0.069 & 0.459 \\
vesicles & 0 & \textbf{ADMM} & kappa=0.01, mu=0.5 & 0.021 & 42.9 & 8 & 1 & 5.3e-06 \\
vesicles & 0.02 & \textbf{MCMC} & denoiser=TV Bregman, sigma=0.02, lambda\_rr=38.863, iter=300 & 0.512 & 29.0 & 61 & 1 & 0.609 \\
vesicles & 0.05 & \textbf{PNP} & denoiser=Gaussian, sigma=25, lambda\_kz=12, iter=40 & 0.732 & 27.5 & 98 & 0 & 0.435 \\
\bottomrule
\end{tabular}\end{adjustbox}\end{minipage}\end{center}
}%
  {\emph{[pending the run]}}
\IfFileExists{tables/cross_comparison_deconv.tex}{\begin{center}\begin{adjustbox}{max width=\linewidth}
\begin{tabular}{llllllll}
\toprule
dataset & noise & best (by NMSE) & params & NMSE & PSNR & SSIM & chi2\_ratio \\
\midrule
img\_001 & 0 & \textbf{ADAM} & reg=none, lambda=0 & 0.0191 & 24.4 & 0.694 & 1.96e-07 \\
img\_001 & 0.02 & \textbf{MCMC} & denoiser=TV Bregman, sigma=0.03 & 0.0245 & 23.4 & 0.596 & 0.98 \\
img\_001 & 0.05 & \textbf{ADAM} & reg=tikhonov, lambda=0.1 & 0.0371 & 21.5 & 0.511 & 1.09 \\
\bottomrule
\end{tabular}\end{adjustbox}\end{center}
}{}
On clean data ADMM (sparsity) wins on all three structures; as noise rises the plug-and-play family
and MCMC take over.

\paragraph{Which denoiser.}
There is \textbf{no universal best denoiser} --- it depends on the object, the noise, and the
algorithm:
\begin{table}[H]\centering
\caption{Winning denoiser (best NMSE) per structure, noise, and plug-and-play algorithm.}
\label{tab:denoisers}
\IfFileExists{tables/denoiser_winners.tex}{\begin{center}\begin{adjustbox}{max width=\linewidth}
\begin{tabular}{lllll}
\toprule
structure & noise & PnP & ADMM-PnP & MCMC \\
\midrule
cell & 0 & DCT (0.191) & Gaussian (0.157) & Gaussian (0.227) \\
fibres & 0 & Gaussian (0.188) & DCT (0.23) & Gaussian (0.115) \\
vesicles & 0 & TV Bregman (0.303) & Gaussian (0.311) & Gaussian (0.166) \\
cell & 0.02 & TV Bregman (0.262) & Gaussian (0.177) & TV Bregman (0.235) \\
fibres & 0.02 & Gaussian (0.385) & DCT (0.568) & DCT (0.511) \\
vesicles & 0.02 & Gaussian (0.522) & DCT (0.603) & DCT (0.588) \\
cell & 0.05 & TV Bregman (0.437) & Gaussian (0.307) & TV Bregman (0.453) \\
fibres & 0.05 & Gaussian (0.653) & TV Bregman (0.862) & DCT (0.856) \\
vesicles & 0.05 & Gaussian (0.732) & TV Bregman (0.87) & DCT (0.889) \\
\bottomrule
\end{tabular}\end{adjustbox}\end{center}
}%
  {\emph{[pending the run]}}
\end{table}
Two facts stand out. \emph{One firm algorithm rule}: TV~Bregman collapses inside the dual-variable
ADMM-PnP on every structure, while it is competitive-to-best in the dual-free PnP and MCMC --- the
dual, more than a single ``best denoiser'', flips the ranking. \emph{DCT is rehabilitated}: long
thought unusable on MA-TIRF, it is often second, sometimes first, once fed at the right scale
($\div\text{peak}$). Bilateral and Wiener, consistently worst, were dropped.

\paragraph{Cost and robustness.}
Runtime, peak memory and rejection rate per solver --- including the PnP family's memory pressure on
the full \texttt{esoubies} volume (every solver is run on it):
\IfFileExists{tables/rejects_and_cost.tex}{\begin{center}\begin{minipage}{\linewidth}\centering
\captionof{table}{Cost and robustness per solver --- runs, rejection rate, failures, median runtime, peak memory and top findings.}
\begin{adjustbox}{max width=\linewidth}
\begin{tabular}{lllllll}
\toprule
solver & runs & rejected & failed/timeout & median s & max mem (MB) & top findings \\
\midrule
ADAM & 45 & 4 (8\%) & 0 & 17.2 & 2176 & no-better×72, depth-shift×18, collapsed×4 \\
ADMM & 19 & 1 (5\%) & 0 & 1.7 & 2204 & no-better×20, depth-shift×10, collapsed×1 \\
ADMM-PnP & 37 & 3 (8\%) & 1 & 14.7 & 10210 & no-better×61, depth-shift×18, collapsed×2 \\
MCMC & 39 & 0 (0\%) & 0 & 9.7 & 2236 & no-better×47, depth-shift×17 \\
PNP & 37 & 0 (0\%) & 0 & 3.7 & 3591 & no-better×40, depth-shift×18 \\
PPXA & 2 & 0 (0\%) & 0 & 291.7 & 1656 & no-better×3, depth-shift×1 \\
\bottomrule
\end{tabular}\end{adjustbox}\end{minipage}\end{center}
}{}

\paragraph{The pattern.}
No single algorithm wins everywhere: the report's thesis is a map, not a champion. Match the prior
to the structure and the noise, and match the denoiser to \emph{both} the object and the
algorithm's mechanism (dual or not).
